\section{Future work}

Several questions follow directly from these experiments.

First, what quantity is being maintained. The data are consistent with conservation of some measure of information flow or representational capacity across heads, but I have not yet measured these directly. Future work will track mutual information, effective rank, gradient norms, and attention budgets over training to identify which quantity remains stable under compensation.

Second, how general is the effect. Parity is a clean Zone 2 task but narrow. The same kill test should be applied to finite state automata, small algorithmic reasoning tasks, and eventually to layer 0 suppressors in full scale language models. If homeostatic compensation is a general property of multi head attention rather than a quirk of this setup it should appear in those settings as well.

Third, how symmetric is compensation. The present experiments focus on $\omega = 0.5$. Prior omega sweeps show symmetric acceleration at $\omega = 0.5$ and $\omega = 1.5$. Running parallel kill tests at $\omega = 1.5$ and freezing different numbers of heads would reveal whether compensation capacity is symmetric and distributed or asymmetric and concentrated.

Finally, can I harness compensation to steer learning. If compensation both accelerates escape from bad plateaus and favours generalizing circuits I may be able to design schedules that deliberately perturb and then restore suppressor strength to shape developmental windows. Doing so safely will require a clearer picture of the equilibrium manifold and of the failure modes revealed by seeds like seed 2.
